\chapter{Ampliación: estabilidad dinámica y linealización}\label{ch:dynamic}
\objetivos{Construir un modelo perturbado; interpretar autovalores; distinguir estabilidad estática y dinámica; reconocer límites del modelo lineal.}

\caja{Contenido de ampliación}{La guía ULE consultada incluye estabilidad estática longitudinal y lateral-direccional e introducción a derivadas. Este capítulo añade fundamentos de estabilidad dinámica para conectar esas derivadas con la simulación. No se presenta como transcripción de un bloque oficial adicional.}

\section{Linealizar alrededor de un movimiento}
Sea $\dot{\vect x}=\vect f(\vect x,\vect u)$ y un equilibrio $\vect f(\vect x_0,\vect u_0)=0$. Con $\Delta\vect x=\vect x-\vect x_0$ y $\Delta\vect u=\vect u-\vect u_0$, el desarrollo de Taylor da
\begin{equation}
\Delta\dot{\vect x}=A\Delta\vect x+B\Delta\vect u+O(\|\Delta\vect x,\Delta\vect u\|^2),
\end{equation}
con $A=\partial\vect f/\partial\vect x$ y $B=\partial\vect f/\partial\vect u$ evaluadas en el trim. Para un movimiento estacionario con coordenadas que evolucionan, como posición o guiñada en viraje, se usan variables apropiadas respecto a la trayectoria de referencia o un modelo lineal variable en el tiempo.

Cada fila de $A$ tiene unidades compatibles con la derivada temporal del estado correspondiente. Un cambio de unidades de estado modifica numéricamente $A$ por una transformación de semejanza y conserva autovalores si se realiza coherentemente. Copiar una matriz de otra referencia sin sus estados y unidades no define un modelo.

\section{Autovalores y respuesta libre}
Para $\Delta\dot{\vect x}=A\Delta\vect x$, la solución es $\Delta\vect x(t)=e^{At}\Delta\vect x(0)$. Un autovalor $\lambda=\sigma+i\omega_d$ genera una contribución con envolvente $e^{\sigma t}$ y oscilación de frecuencia angular $|\omega_d|$.

Si todas las partes reales son negativas, el equilibrio lineal es asintóticamente estable. Una parte real positiva implica inestabilidad lineal. Partes reales nulas requieren examinar estructura de Jordan y términos no lineales; no basta con clasificarlas como estabilidad garantizada. En aviones pueden aparecer integradores de posición y rumbo que no significan inestabilidad de actitud.

Para un par complejo amortiguado,
\begin{equation}
\omega_n=\sqrt{\sigma^2+\omega_d^2},\qquad
\zeta=-\sigma/\omega_n,\qquad T_d=2\pi/|\omega_d|.
\end{equation}
El tiempo de reducción a la mitad de la envolvente es $\ln2/(-\sigma)$ si $\sigma<0$. El periodo de oscilación y el tiempo de amortiguamiento describen propiedades diferentes.

\ilustra{modos-dinamicos}{Dos respuestas sintéticas con la misma frecuencia de oscilación y signos opuestos en la tasa exponencial. Una tendencia restauradora puede coexistir con una oscilación creciente si falta amortiguamiento.}

\section{Estabilidad estática no implica amortiguamiento}
Un modelo de cabeceo elemental $I_y\ddot\alpha+c\dot\alpha+k\alpha=0$ tiene rigidez restauradora si $k>0$. Si $c<0$, aporta energía a la perturbación y la oscilación puede crecer. La condición estática está relacionada con rigidez; la dinámica requiere además amortiguamiento y acoplamientos.

Este ejemplo no convierte $\alpha$ en una coordenada rígida exacta de cabeceo: la incidencia depende de actitud y trayectoria. Sirve para mostrar la diferencia lógica. Un modelo de avión se construye con las ecuaciones linealizadas de fuerzas y momentos.

\section{Modelo longitudinal con derivadas de aceleración}
En ejes de estabilidad cerca de vuelo simétrico se puede elegir $\vect x=(u,w,q,\theta)^T$. Una forma descriptor conserva términos como $Z_{\dot w}$ y $M_{\dot w}$:
\begin{equation}
E_d\dot{\vect x}=A_d\vect x+B_d\vect u.
\end{equation}
Si $E_d$ es invertible, se calcula $A=E_d^{-1}A_d$. No se divide sólo una fila por la masa si el término de aceleración aparece acoplado también en la ecuación de momento. Es preferible resolver el sistema lineal que formar una inversa explícita en un algoritmo numérico.

Las derivadas dimensionales $X_u,Z_w,M_q$ incluyen las normalizaciones de fuerza y momento. Si se parte de coeficientes, se aplican regla de la cadena y factores $q_\infty S$, $\bar c/(2V)$ y relaciones entre $w$ y $\alpha$. La densidad, velocidad y masa del trim forman parte del resultado.

\section{Modos longitudinales}
El periodo corto suele involucrar principalmente incidencia y velocidad angular de cabeceo, con variación de velocidad relativamente pequeña en su escala temporal. La fugoide intercambia principalmente energía cinética y potencial y suele ser más lenta. La participación exacta de cada estado se observa en autovectores y respuesta, no sólo en el nombre del par de autovalores.

Una aproximación energética sin pérdidas de la fugoide da una escala $\omega\simeq\sqrt2g_0/V_0$. No predice amortiguamiento y depende de hipótesis fuertes sobre sustentación y energía. Debe usarse como orden de magnitud, no como reemplazo de la matriz completa.

\caso{Ejemplo original de periodo corto sintético}{Para estados $(\alpha,q)$ en radianes y radianes por segundo se toma
\[
A=\begin{bmatrix}-0.8&1.0\\-4.0&-1.2\end{bmatrix}.
\]
El polinomio característico es $\lambda^2+2\lambda+4.96=0$, con $\lambda=-1\pm1.990i\SIu{s^{-1}}$. Resultan $\omega_n=2.227\SIu{rad\,s^{-1}}$, $\zeta=0.449$ y periodo $3.16\SIu s$. Se comprueban traza y determinante frente a suma y producto de autovalores. Es una matriz docente y no un modelo identificado del T-6 ni de otra aeronave.}

\section{Modos laterales}
En una reducción lateral se usan habitualmente $(\beta,p,r,\phi)$. El modo de alabeo suele ser una raíz real rápida asociada a amortiguamiento de $p$. El balanceo holandés es un par oscilatorio acoplado de guiñada y resbalamiento, con alabeo participante. El modo espiral es habitualmente una raíz real lenta cuyo signo depende de la competencia entre tendencias direccionales, efecto diedro y acoplamientos.

Un avión puede tener $C_{n\beta}>0$ y $C_{\ell\beta}<0$ y aun así presentar una divergencia espiral lenta. El diagnóstico exige la matriz completa en el trim. Los productos de inercia introducen acoplamientos entre aceleraciones de alabeo y guiñada y deben conservarse o descartarse justificadamente.

\section{Extracción numérica de una matriz}
La diferencia centrada de la función de estado permite estimar cada columna de $A$. Se deben restaurar estado y mandos exactamente en cada evaluación, evitar que el motor arrastre estados internos de una perturbación a otra y repetir con varios tamaños de paso. Un simulador con memoria de actuadores o motor requiere incluir esos estados o congelarlos de manera documentada.

Una validación interna compara $e^{At}\Delta\vect x_0$ con una simulación no lineal para perturbaciones pequeñas. La discrepancia debe reducirse al disminuir amplitud dentro del intervalo donde el ruido numérico no domina. Coincidencia local no demuestra fidelidad experimental del modelo.

\section{Control y límites de la estabilidad lineal}
Un piloto automático modifica la dinámica cerrada. Con realimentación $\Delta\vect u=-K\Delta\vect x$, la matriz pasa a $A-BK$ dentro del modelo lineal. Saturación, retrasos, límites de actuador y cambios de condición pueden invalidar esa representación. Un control que mantiene altura en una simulación no convierte el avión libre en estable ni certifica el sistema.

\section{Construcción explícita del modelo longitudinal}
Se adoptan ejes de estabilidad alineados con la velocidad del trim, vuelo simétrico a velocidad $U_0$ y cabeceo $\theta_0$. Las perturbaciones son $(u,w,q,\theta)$. Las fuerzas no gravitatorias se linealizan como $\Delta X=X_uu+X_ww+X_qq+\cdots$ y análogamente Z y M. Las derivadas dimensionales incluyen su condición de evaluación.

Si se desprecia $X_q$ para esta forma reducida y se conservan derivadas de aceleración $Z_{\dot w}$ y $M_{\dot w}$,
\begin{align}
m\dot u&=X_uu+X_ww-mg_0\cos\theta_0\,\theta+\Delta X_c,\\
(m-Z_{\dot w})\dot w&=Z_uu+Z_ww+(Z_q+mU_0)q-mg_0\sin\theta_0\,\theta+\Delta Z_c,\\
I_{yy}\dot q-M_{\dot w}\dot w&=M_uu+M_ww+M_qq+\Delta M_c,\\
\dot\theta&=q.
\end{align}
La forma descriptor utiliza
\begin{equation}
E_d=\begin{bmatrix}m&0&0&0\\0&m-Z_{\dot w}&0&0\\0&-M_{\dot w}&I_{yy}&0\\0&0&0&1\end{bmatrix}.
\end{equation}
El coeficiente $m-Z_{\dot w}$ no es una masa añadida arbitraria: reúne una dependencia de fuerza con aceleración del estado. Su signo y magnitud proceden del modelo aerodinámico. Si la matriz descriptor es singular o mal condicionada, el paso a forma estándar debe revisarse en lugar de dividir sin diagnóstico.

\subsection{Unidades de derivadas}
$X_u$ tiene unidad N/(m/s), $M_q$ tiene N m/(rad/s) y $Z_{\dot w}$ tiene kg al tratar radianes como adimensionales. Una derivada de coeficiente respecto a $q\bar c/(2V)$ se convierte en $M_q=q_\infty S\bar c\,C_{mq}\bar c/(2V)$. Si se olvida el segundo factor de cuerda, la dimensión del momento derivado es incorrecta.

La relación $\alpha\simeq w/U_0$ permite convertir $Z_\alpha$ en $Z_w\simeq Z_\alpha/U_0$ bajo el trim elegido. Derivar $q_\infty$ respecto a u también aporta términos: una derivada de fuerza no es siempre sólo $qS$ por la derivada del coeficiente. La regla de la cadena debe aplicarse a toda la fuerza.

\section{Una reducción de periodo corto}
Si u cambia poco en la escala rápida, se conservan principalmente $\alpha\simeq w/U_0$ y q. En una forma genérica,
\begin{equation}
\begin{bmatrix}\dot\alpha\\\dot q\end{bmatrix}
=\begin{bmatrix}a_{11}&a_{12}\\a_{21}&a_{22}\end{bmatrix}
\begin{bmatrix}\alpha\\q\end{bmatrix}.
\end{equation}
Su polinomio es $\lambda^2-\operatorname{tr}(A_{sp})\lambda+\det(A_{sp})=0$. Para estabilidad asintótica de un sistema real de segundo orden, traza negativa y determinante positivo son condiciones necesarias y suficientes. La presencia de oscilación requiere además $\operatorname{tr}(A_{sp})^2<4\det(A_{sp})$.

Un $C_{m\alpha}$ negativo suele contribuir a rigidez restauradora, pero $a_{21}$ incluye normalizaciones y posibles derivadas de aceleración. La traza contiene amortiguamientos; por ello la estabilidad estática por sí sola no fija el signo de las partes reales. La reducción debe compararse con el sistema de cuatro estados para verificar la separación de escalas.

\section{Derivación energética de la escala de fugoide}
Supóngase una perturbación lenta en la que se conserva aproximadamente energía específica $V^2/2+g_0h$ y CL. Al linealizar alrededor de $U_0,h_0$,
\begin{equation}
U_0u+g_0\Delta h\simeq0,\qquad u\simeq-g_0\Delta h/U_0.
\end{equation}
La variación de sustentación por velocidad es $\Delta L\simeq\rho SC_LU_0u$. La aceleración vertical aproximada es $m\Delta\ddot h=\Delta L$. Con $\rho SC_LU_0^2/2=mg_0$ se obtiene
\begin{equation}
\Delta\ddot h+\frac{2g_0^2}{U_0^2}\Delta h\simeq0.
\end{equation}
La frecuencia $\sqrt2g_0/U_0$ es una escala, no una predicción completa: se omitieron resistencia, consumo, variación de CL y control. El modelo no predice amortiguamiento porque se supuso conservación de energía. Incorporarlo exige recuperar fuerzas disipativas y acoplamientos.

\section{Modelo lateral con producto de inercia}
Para perturbaciones $(v,p,r,\phi)$ en vuelo simétrico de cabeceo $\theta_0$,
\begin{align}
m\dot v&=Y_vv+Y_pp+(Y_r-mU_0)r+mg_0\cos\theta_0\phi+\Delta Y_c,\\
I_{xx}\dot p-I_{xz}\dot r&=L_vv+L_pp+L_rr+\Delta L_c,\\
-I_{xz}\dot p+I_{zz}\dot r&=N_vv+N_pp+N_rr+\Delta N_c,\\
\dot\phi&=p+r\tan\theta_0.
\end{align}
Se ha linealizado alrededor de tasas de trim nulas; en un trim de viraje aparecen términos adicionales. La ecuación lateral conserva $-U_0r$ por la aceleración de transporte. El bloque de inercias se resuelve antes de ensamblar A. Copiar derivadas $L_p/I_{xx}$ y $N_r/I_{zz}$ como si no hubiera producto de inercia elimina un acoplamiento real del modelo.

\subsection{Amortiguamiento de alabeo como límite}
Si v y r varían poco y se desprecia Ixz, $\dot p\simeq(L_p/I_{xx})p$. Con $L_p<0$, la tasa decae con constante temporal $-I_{xx}/L_p$. Esta aproximación explica una raíz real rápida. No describe el ángulo final de alabeo sin integrar $\dot\phi=p$ ni incluye el modo espiral lento.

\section{Autovectores y observación de modos}
Los autovalores indican tasas y frecuencias; los autovectores indican combinaciones de estados que participan. Un modo puede ser poco visible en una variable y evidente en otra. Además, la amplitud observada depende de la proyección de la perturbación inicial sobre los modos y de las salidas medidas. No encontrar una oscilación en velocidad no demuestra ausencia del periodo corto si éste aparece principalmente en incidencia y q.

Los autovectores cambian con el escalado de estados. Para comparar participación se normalizan con escalas físicas o se interpretan relaciones concretas y fase. Una tabla de magnitudes sin unidades puede dar una impresión engañosa de qué variable domina.

\section{Límite de validez de una linealización}
La aproximación descarta términos de segundo orden. Para comprobarla se compara con el sistema no lineal a amplitudes decrecientes. Si se reduce la perturbación a la mitad, una discrepancia dominada por términos cuadráticos debería disminuir aproximadamente más deprisa que la propia perturbación, antes de que dominen redondeo y tolerancias. Cerca de una discontinuidad de tabla, pérdida, saturación o singularidad, puede no existir una región lineal útil con el paso elegido.

La respuesta del sistema con piloto automático se interpreta con estados de actuadores e integradores del control. Una matriz de avión libre no predice exactamente esa trayectoria cerrada. En el experimento del 747 se utiliza control docente de seguimiento; por ello no se infiere estabilidad libre de la aeronave a partir de que conserve aproximadamente la altura.


\section{Taller resuelto: Respuesta lineal por dos integradores conceptuales}

\noindent\textbf{Nivel:} Intermedio.

\apartado{1. Enunciado}{Para la matriz sintética A=[[-0.7,1],[-3.8,-1.1]] y perturbación inicial (alpha,q)=(0.04,0), calcule autovalores y estado a 12 s. Compare exponencial matricial y una integración de ODE.}

\apartado{2. Datos}{alpha en rad, q en rad/s; matriz no identificada de una aeronave.}

\apartado{3. Incógnitas}{Las magnitudes y comprobaciones solicitadas en el caso.}

\apartado{4. Sistema de referencia}{Estados longitudinales perturbados; tiempo en segundos.}

\apartado{5. Hipótesis físicas}{Sistema lineal autónomo, sin mando y sin saturación.}

\apartado{6. Modelo matemático}{Respuesta lineal por dos integradores conceptuales}

\apartado{7. Ecuaciones simbólicas}{Se parte de los balances y definiciones siguientes.}

\begin{equation}\dot{\vect x}=A\vect x\end{equation}

\begin{equation}\vect x(t)=e^{At}\vect x(0)\end{equation}

\apartado{8. Desarrollo algebraico}{Las operaciones se realizan en este orden, conservando las condiciones del modelo.}

\begin{enumerate}

\item Se comprueba que ambas partes reales son negativas.

\item Se evalúa la exponencial de matriz y se integra la ODE con las mismas condiciones.

\item Se compara el estado completo y no sólo un cruce de alpha por cero.

\item La diferencia numérica mide consistencia de estos algoritmos sobre el modelo lineal.

\end{enumerate}

\apartado{9. Sustitución numérica}{Se conservan los resultados intermedios sin redondear y se muestran todas las operaciones escalares disponibles.}

\textbf{Exponencial desarrollada para esta matriz.} El polinomio característico es $\lambda^2+1.8\lambda+4.57=0$, de donde $\lambda=-0.9\pm i\sqrt{3.76}$. Escribir $A=-0.9I+B$ da $B=\bigl(\begin{smallmatrix}0.2&1\\-3.8&-0.2\end{smallmatrix}\bigr)$ y verificar por multiplicación $B^2=-3.76I$. Agrupar términos pares e impares de la serie exponencial produce:\[\fitmath{e^{At}=e^{-0.9t}\left[I\cos(\sqrt{3.76}t)+\frac{B}{\sqrt{3.76}}\sin(\sqrt{3.76}t)\right]}\]\[\fitmath{\alpha(t)=0.04e^{-0.9t}\left[\cos(\sqrt{3.76}t)+\frac{0.2}{\sqrt{3.76}}\sin(\sqrt{3.76}t)\right]}\]\[\fitmath{q(t)=-\frac{0.152}{\sqrt{3.76}}e^{-0.9t}\sin(\sqrt{3.76}t)}\]Sustituir $t=12$ s da $\alpha=-3.163128\times10^{-7}\,\mathrm{rad}$ y $q=1.530889\times10^{-6}\,\mathrm{s^{-1}}$. Para contrastar con integración, se usa el mismo estado inicial y se controla cada componente por separado; una diferencia entre algoritmos no es validación experimental.





\apartado{10. Resultados con unidades}{Los resultados pertenecen al modelo y al intervalo de datos declarados.}

\begin{center}\begin{tabular}{lr}\toprule Magnitud&Resultado\\\midrule

Parte real modal & $-0.9\,\mathrm{s^{-1}}$\\

Parte imaginaria modal & $1.939072\,\mathrm{s^{-1}}$\\

Alpha a 12 s & $-3.163128\times10^{-7}\,\mathrm{rad}$\\

q a 12 s & $1.530889\times10^{-6}\,\mathrm{s^{-1}}$\\

Error de alpha entre métodos & $4.039651\times10^{-14}\,\mathrm{rad}$\\

Error de q entre métodos & $4.606294\times10^{-14}\,\mathrm{s^{-1}}$\\

\bottomrule\end{tabular}\end{center}

\apartado{11. Interpretación}{El estado final pequeño corresponde a una oscilación amortiguada; el resultado pertenece a la matriz sintética.}

\apartado{12. Comprobación independiente y alcance}{La comparación expm frente a $solve_{ivp}$ se ejecuta en el programa; se conservan tolerancias e intervalo.}

\apartado{13. Limitaciones}{La norma de error agrupa dos estados de unidades distintas sólo como diagnóstico numérico; cada componente se revisa además individualmente.}

\par\smallskip\noindent{\footnotesize\textit{Datos y solución originales. Las sustituciones registradas se reproducen con SymPy y Studium; eso no establece validación experimental ni revisión experta.}}


\repaso{¿Basta con ver que una perturbación vuelve a cero una vez?}{No. Una oscilación creciente cruza cero repetidamente. Hay que examinar la envolvente, varias variables, tiempo suficiente y condiciones de mando; los autovalores aportan una descripción local complementaria.}
\resumen{La linealización vincula derivadas con modos de movimiento. Los signos estáticos no fijan por sí solos las partes reales de los autovalores. Cada matriz necesita estados, unidades y trim explícitos, y la comparación con el modelo no lineal conserva un alcance local.}
\fuentes{How, MIT 16.333, lección 4 (linealización), lección 6, pp. 5--1 a 5--3 (forma descriptor), lección 7, pp. 6--1 a 6--4 (aproximaciones), y lección 8, pp. 7--1 a 7--3 (dinámica lateral). Ejemplo matricial y figuras originales.}
